\medskip

L'espace est rapporté à un repère orthonormé \Oijk.

On considère les points 
\begin{center}A$(1~;~0~;~-1)$,\quad  B$(3~;~-1~;~2)$, \quad C$(2~;~-2~;~ -1)$ \quad et \quad D$(4~;~-1~;~-2)$.\end{center}

On note $\Delta$ la droite de représentation paramétrique 
\[\left\{\begin{array}{l c l}
x &=&\phantom{-}0\\
y &=&\phantom{-}2 + t\\
z &=& -1 + t
\end{array}\right.\:,\:\: \text{avec } t\in \R.\]

\begin{enumerate}
\item 
	\begin{enumerate}
		\item 

On a $\vect{\text{AB}}\begin{pmatrix}2\\-1\\3\end{pmatrix}$ et $\vect{\text{AC}}\begin{pmatrix}1\\-2\\0\end{pmatrix}$ : ces deux vecteurs ne sont pas colinéaires (en comparant les premières coordonnées, s'il l'étaient on aurait $\vect{\text{AB}} = 2\vect{\text{AC}}$, ce qui est faux), donc les points A, B et C ne sont pas alignés et définissent donc un plan unique $\mathcal{P}$.
		\item 
On a $\vect{\text{CD}}\begin{pmatrix}2\\1\\-1 \end{pmatrix}$, puis 

$\bullet~~$ $\vect{\text{CD}} \cdot \vect{\text{AB}} = 4 - 1 - 3 = 0$ ;

$\bullet~~$ $\vect{\text{CD}} \cdot \vect{\text{AC}} = 2 - 2 + 0 = 0$.

Le vecteur $\vect{\text{CD}}$ est donc orthogonal à deux vecteurs non colinéaires du plan $\mathcal{P}$, c'est donc un vecteur normal à ce plan.


Puisque $\vect{\text{CD}} \cdot \vect{\text{AC}}  = 0$ les droites (CD) et (AC) sont perpendiculaires en C, donc C est le projeté orthogonal de D sur le plan $\mathcal{P}$.
		\item 
		Puisque $\vect{\text{CD}}$ est un vecteur normal au plan $\mathcal{P}$, on sait que ses cordonnées sont dans l'ordre les coefficients $a, \: b,\:c$ de $x,\:y,\:z$, soit :
		
$M(x~;~y~;~z) \in \mathcal{P} \iff 2x + 1y - 1z + d = 0$.

En particulier C$(2~;~-2~;~ -1) \in \mathcal{P} \iff 2\times 2  + 1\times (- 2)  - 1\times (- 1) + d = 0 \iff 4 - 2 + 1 + d = 0 \iff d = - 3$.

Finalement : $M(x~;~y~;~z) \in \mathcal{P} \iff 2x + y - z - 3 = 0$.
	\end{enumerate}
\item 
	\begin{enumerate}
		\item 
On a grâce aux coordonnées du vecteur $\vect{\text{CD}}$, \:CD$^2 = 4 + 1 + 1 = 6$, d'où CD $ = \sqrt 6$.
		\item 
Puisque C est le projeté orthogonal de D sur le plan $\mathcal{P}$, ce point est celui qui est à la plus courte distance du point D, soit $\sqrt 6$ : il n'existe donc pas d'autre point de $\mathcal{P}$ situé à cette distance $\sqrt 6$ de D;

Ou encore : la sphère de centre D et de rayon $\sqrt 6$ est tangente en C au plan $\mathcal{P}$, donc quel que soit M dans $\mathcal{P}$, le triangle DCM est rectangle en C, d'hypoténuse [DM] et l'on sait qu'alors DM > DC $= \sqrt 6$ : tout point M du plan aurez que C est donc à une distance supérieure à $\sqrt 6$ de D.
	\end{enumerate}
\item 
	\begin{enumerate}
		\item 
Si $M(x~;~y~;~z)$ est commun à $\Delta$ et à $\mathcal{P}$ ses coordonnées vérifient le système :

$\left\{\begin{array}{l c l}
x &=&\phantom{-}0\\
y &=&\phantom{-}2 + t\\
z &=& -1 + t\\
2x + y - z - 3 &=& 0
\end{array}\right.\:,\:\: \text{avec } t\in \R \Rightarrow 2 \times 0  + 2 + t  - (- 1 + t) - 3 = 0 \iff 0 = 0$ : ceci signifie que tout point de $\Delta$ appartient à $\mathcal{P}$ donc que $\Delta$ est incluse dans le plan $\mathcal{P}$.

Soit H le projeté orthogonal du point D sur la droite $\Delta$.
		\item 
Un vecteur directeur de la droite $\Delta$ est $\delta\begin{pmatrix}0\\1\\1\end{pmatrix}$ et avec H$(0~;~2 + t~;~- 1 + t)$, on a $\vect{\text{DH}}\begin{pmatrix}4\\3 + t\\1 + t\end{pmatrix}$.

Or (DH) $\perp \Delta \Rightarrow \delta \cdot \vect{\text{DH}} = 0 \iff 0 + 3 + t + 1 + t = 0 \iff 4 + 2t = 0 \iff 2t = - 4 \iff t = - 2$.

On a donc H$(0~;~0~;~- 3) \in \Delta$.
		\item 
On a donc $\vect{\text{DH}}\begin{pmatrix}4\\3 - 2\\1 - 2\end{pmatrix}$, soit $\vect{\text{DH}}\begin{pmatrix}4\\1\\- 1\end{pmatrix}$ ; il en résulte que DH$^2 = 4^2 +1^2 + (- 1)^2 = 18$.
		
Finalement DH $= \sqrt{18} = \sqrt{9 \times 2} = \sqrt 9 \times \sqrt 2 = 3\sqrt 2$.
	\end{enumerate}
\end{enumerate}


\bigskip
